% !TeX program = xelatex
\documentclass[12pt,a4paper]{article}
\usepackage[margin=2cm]{geometry}
\usepackage{fontspec}
\usepackage{amsmath,amssymb,enumitem}
\setmainfont{Times New Roman}
\setlength{\parindent}{0pt}\setlength{\parskip}{6pt}
\begin{document}
\begin{center}{\Large\bfseries Đáp án Toán vào lớp 10 Hà Nội năm 2015\\Phần Đại số}\end{center}
\textbf{Câu 33.}
\begin{enumerate}[label=\alph*)]
\item Thay $x=9$:
\[
P=\frac{9+3}{\sqrt9-2}=\frac{12}{3-2}=12.
\]
\item Với $x>0$, $x\ne4$:
\begin{align*}
Q&=\frac{\sqrt{x}-1}{\sqrt{x}+2}+\frac{5\sqrt{x}-2}{x-4}\\
&=\frac{(\sqrt{x}-1)(\sqrt{x}-2)+5\sqrt{x}-2}{x-4}\\
&=\frac{x+2\sqrt{x}}{(\sqrt{x}-2)(\sqrt{x}+2)}
=\frac{\sqrt{x}}{\sqrt{x}-2}.
\end{align*}
\item Từ câu b),
\[
\frac{P}{Q}=\frac{x+3}{\sqrt{x}}=\sqrt{x}+\frac{3}{\sqrt{x}}\ge2\sqrt3,
\]
do bất đẳng thức AM--GM. Dấu bằng xảy ra khi $\sqrt{x}=3/\sqrt{x}$, tức $x=3$ (thỏa điều kiện). Vậy giá trị nhỏ nhất là $2\sqrt3$, đạt tại $x=3$.
\end{enumerate}

\textbf{Câu 34.} Gọi $x$ (km/giờ) là vận tốc tàu khi nước yên lặng; $x>2$. Thời gian ngược dòng và xuôi dòng lần lượt là $60/(x-2)$ và $48/(x+2)$ giờ. Theo đề bài,
\[
\frac{60}{x-2}-\frac{48}{x+2}=1
\iff x^2-12x-220=0
\iff x=22\text{ hoặc }x=-10.
\]
Do $x>2$, nhận $x=22$. Vận tốc cần tìm là $22$ km/giờ.

\textbf{Câu 35.} Điều kiện $x\ge-1$. Đặt $a=x+y$, $b=\sqrt{x+1}\ (b\ge0)$. Hệ trở thành
\[
\begin{cases}2a+b=4,\\a-3b=-5,\end{cases}
\quad\Rightarrow\quad a=1,\ b=2.
\]
Suy ra $x+1=4$, nên $x=3$, và $x+y=1$ cho $y=-2$. Nghiệm của hệ là $(3;-2)$.

\textbf{Câu 36.}
\begin{enumerate}[label=\alph*)]
\item $\Delta=(m+5)^2-4(3m+6)=(m-1)^2\ge0$ với mọi $m\in\mathbb R$. Vậy phương trình luôn có nghiệm thực.
\item Phân tích đa thức:
\[
x^2-(m+5)x+3m+6=(x-3)(x-m-2),
\]
nên hai nghiệm là $3$ và $m+2$. Vì là độ dài hai cạnh góc vuông, cần $m+2>0$ và
\[
3^2+(m+2)^2=5^2.
\]
Suy ra $m+2=4$ (do dương), vậy $m=2$.
\end{enumerate}

\textbf{Câu 38.} Từ $a^2+b^2=4$,
\[
2ab=(a+b)^2-4,\qquad 2M=\frac{(a+b)^2-4}{a+b+2}=a+b-2.
\]
Mặt khác $a+b\le\sqrt{2(a^2+b^2)}=2\sqrt2$. Do đó $M\le\sqrt2-1$. Dấu bằng khi $a=b=\sqrt2$, thỏa điều kiện. Vậy $\max M=\sqrt2-1$.
\end{document}
